\vspace{-0.2cm}
\bluenote{
% 基于软件和库，代码量
We implement our FedFT prototype based on the open-source FedPETuning framework \cite{zhang2022federated}, extending its functionality with approximately 2.1K lines of custom code. 
The prototype is designed to support heterogeneous computing platforms, with a specific focus on commercial NVIDIA Jetson devices \cite{mittal2019survey}, including Jetson TX2, Jetson NX, and Jetson AGX.
The software platform is built based on Docker Swarm \cite{merkel2014docker, naik2016building}, a distributed software development kit that helps build distributed systems with the ability to monitor the status of each device.
For computing, we utilize PyTorch \cite{paszke2019pytorch} to facilitate the implementation of model fine-tuning on devices, ensuring platform independence while allowing platform-specific backend acceleration. 
To accelerate the on-device fine-tuning, we utilize NVIDIA-provided development packages\footnote{https://forums.developer.nvidia.com/t/pytorch-for-jetson/72048} to take full advantage of the underlying hardware capabilities.
For communication, we adopt MPI (Message Passing Interface) \cite{gabriel2004open}, which includes a collection of sending and receiving functions, \eg, \textit{comm.send(data, dest, tag)}/\textit{comm.recv(sour, tag)}, to streamline communication between the PS and devices.
The prototype provides simple APIs to abstract away the problem of federated fine-tuning on heterogeneous computing platforms.
For example, in the PS, we use docker API \textit{docker stack deploy} to deploy the fine-tuning process on both PS and the heterogeneous device with the customized container for different devices.
The implementation addresses key challenges in supporting heterogeneous devices by developing flexible communication and fine-tuning protocols that can adapt to varying computational resources.
}
\vspace{-0.2cm}



% \rednote{
% The prototype FedFT system comprises two critical modules, \ie, communication and local fine-tuning. 
% The communication module, implemented using Socket protocols, incorporates robust features such as heartbeat detection and efficient data transmission mechanisms. 
% This ensures reliable server-device communication and synchronization in federated learning scenarios.
% The local fine-tuning module is implemented using PyTorch, ensuring platform independence while allowing platform-specific backend acceleration. 
% Critically, the backend libraries are optimized for specific hardware platforms and software versions. 
% For Jetson devices, we utilized NVIDIA-provided development packages to accelerate training processes, taking full advantage of the underlying hardware capabilities.
% PyTorch installation is carefully managed following NVIDIA developer forum instructions \footnote{https://forums.developer.nvidia.com/t/pytorch-for-jetson/72048}, with custom configurations for each Jetson device variant. 
% The implementation addresses key challenges in heterogeneous device support by developing flexible communication and fine-tuning protocols that can adapt to varying computational resources.
% }
